% GENERATED FILE. Edit canonical lessons, then run npm run generate:books.
% canonical-source-hash: fnv1a64:e429848721b8649b
% canonical-lessons: ZH-C136-qingtian, ZH-C136-yintian, ZH-C136-guafeng, ZH-C136-dalei, ZH-C136-liangkuai

\chapter{Because of the Weather}
\label{ch:zh-because-of-the-weather}
\hlchaptermodality{136}

\begin{chapteropening}
\textbf{By the end of this chapter:} \emph{I can give the weather as a reason: sunny, overcast, windy, thundery or cool.}

Complete the last lesson of chapter 136: I can give the weather as a reason: sunny, overcast, windy, thundery or cool.
\end{chapteropening}

\section[qíngtiān]{qíngtiān — sunny day; fine weather}

\label{lesson:ZH-C136-qingtian}



\begin{quote}
Before the new one: say the Mandarin for a match, then the Mandarin for to work out.
\end{quote}

\subsection*{You'll want to know: qíngtiān}
\emph{qíngtiān} — \textquotedblleft{}sunny day; fine weather\textquotedblright{}.

\begin{quote}
\emph{jīntiān shì qíngtiān wǒ xiǎng chūmén} — It's a sunny day, so I feel like going out.
\end{quote}

Its characters come later; for now it is a word for the ear.

\begin{cousinweb}
\emph{qíngtiān} is built from \emph{qíng} 'clear' and \zh{天} 'sky', both inherited from Middle Chinese, \emph{dzjeng} and \emph{then} (Guangyun, in Baxter's transcription); the dictionary glosses the first 'the sky clearing'.
\end{cousinweb}

\begin{grammarlens}[title={What you've built}]
One more word for giving reasons.
\end{grammarlens}

\subsection*{Your turn}
\begin{itemize}
\raggedright
  \item \emph{Say it:} \emph{qíngtiān}
  \item \emph{Say it:} \emph{qíngtiān}, once more
  \item \emph{Say it:} \emph{jīntiān shì qíngtiān wǒ xiǎng chūmén} — It's a sunny day, so I feel like going out.
  \item \emph{Recall:} say the Mandarin for a match, then the Mandarin for to work out, then say \emph{qíngtiān} again
\end{itemize}

\begin{tcolorbox}[breakable,colback=teal!4,colframe=teal!35!black,title={Before you move on}]
What does qíngtiān mean? (\textquotedblleft{}Sunny day; fine weather\textquotedblright{}.) Say it once more.
\end{tcolorbox}
\section[yīntiān]{yīntiān — overcast day}

\label{lesson:ZH-C136-yintian}



\begin{quote}
Before the new one: say the Mandarin for to work out, then the Mandarin for a sunny day.
\end{quote}

\subsection*{You'll want to know: yīntiān}
\emph{yīntiān} — \textquotedblleft{}overcast day\textquotedblright{}.

\begin{quote}
\emph{jīntiān shì yīntiān yǒudiǎnr lěng} — It's overcast today, and a bit cold.
\end{quote}

Its characters come later; for now it is a word for the ear.

\begin{cousinweb}
\emph{yīntiān} is built from \emph{yīn} 'shade, dark' and \zh{天} 'sky', both inherited from Middle Chinese, \emph{'im} and \emph{then} (Guangyun, in Baxter's transcription). The first has the same old form as the \zh{音} of \zh{音乐} but is a different word.
\end{cousinweb}

\begin{grammarlens}[title={What you've built}]
One more word for giving reasons.
\end{grammarlens}

\subsection*{Your turn}
\begin{itemize}
\raggedright
  \item \emph{Say it:} \emph{yīntiān}
  \item \emph{Say it:} \emph{yīntiān}, once more
  \item \emph{Say it:} \emph{jīntiān shì yīntiān yǒudiǎnr lěng} — It's overcast today, and a bit cold.
  \item \emph{Recall:} say the Mandarin for to work out, then the Mandarin for a sunny day, then say \emph{yīntiān} again
\end{itemize}

\begin{tcolorbox}[breakable,colback=teal!4,colframe=teal!35!black,title={Before you move on}]
What does yīntiān mean? (\textquotedblleft{}Overcast day\textquotedblright{}.) Say it once more.
\end{tcolorbox}
\section[guāfēng]{guāfēng — to be windy}

\label{lesson:ZH-C136-guafeng}



\begin{quote}
Before the new one: say the Mandarin for a sunny day, then the Mandarin for an overcast day.
\end{quote}

\subsection*{You'll want to know: guāfēng}
\emph{guāfēng} — \textquotedblleft{}to be windy\textquotedblright{}.

\begin{quote}
\emph{wàibian guāfēng le} — It's got windy outside.
\end{quote}

Its characters come later; for now it is a word for the ear.

\begin{cousinweb}
\emph{guāfēng} is built from \emph{guā} 'to scrape' and \emph{fēng} 'wind', both inherited from Middle Chinese, \emph{kwaet} and \emph{pjuwng} (Guangyun, in Baxter's transcription). For the wind the first was once written with another character, which has no Guangyun reading, so how scraping came into the word is not settled.
\end{cousinweb}

\begin{grammarlens}[title={What you've built}]
One more word for giving reasons.
\end{grammarlens}

\subsection*{Your turn}
\begin{itemize}
\raggedright
  \item \emph{Say it:} \emph{guāfēng}
  \item \emph{Say it:} \emph{guāfēng}, once more
  \item \emph{Say it:} \emph{wàibian guāfēng le} — It's got windy outside.
  \item \emph{Recall:} say the Mandarin for a sunny day, then the Mandarin for an overcast day, then say \emph{guāfēng} again
\end{itemize}

\begin{tcolorbox}[breakable,colback=teal!4,colframe=teal!35!black,title={Before you move on}]
What does guāfēng mean? (\textquotedblleft{}To be windy\textquotedblright{}.) Say it once more.
\end{tcolorbox}
\section[dǎléi]{dǎléi — to thunder}

\label{lesson:ZH-C136-dalei}



\begin{quote}
Before the new one: say the Mandarin for an overcast day, then the Mandarin for to be windy.
\end{quote}

\subsection*{You'll want to know: dǎléi}
\emph{dǎléi} — \textquotedblleft{}to thunder\textquotedblright{}.

\begin{quote}
\emph{dǎléi le wǒmen bù chūmén le} — It's thundering, so we're staying in.
\end{quote}

Its characters come later; for now it is a word for the ear.

\begin{cousinweb}
\emph{dǎléi} is built from \zh{打} 'to strike' and \emph{léi} 'thunder', both inherited from Middle Chinese, \emph{taengX} and \emph{lwoj} (Guangyun, in Baxter's transcription), though dǎ is not the sound \emph{taengX} would regularly give. The dictionary explains the second, after an older book, as yin and yang pressing on each other.
\end{cousinweb}

\begin{grammarlens}[title={What you've built}]
One more word for giving reasons.
\end{grammarlens}

\subsection*{Your turn}
\begin{itemize}
\raggedright
  \item \emph{Say it:} \emph{dǎléi}
  \item \emph{Say it:} \emph{dǎléi}, once more
  \item \emph{Say it:} \emph{dǎléi le wǒmen bù chūmén le} — It's thundering, so we're staying in.
  \item \emph{Recall:} say the Mandarin for an overcast day, then the Mandarin for to be windy, then say \emph{dǎléi} again
\end{itemize}

\begin{tcolorbox}[breakable,colback=teal!4,colframe=teal!35!black,title={Before you move on}]
What does dǎléi mean? (\textquotedblleft{}To thunder\textquotedblright{}.) Say it once more.
\end{tcolorbox}
\section[liángkuai]{liángkuai — pleasantly cool}

\label{lesson:ZH-C136-liangkuai}



\begin{quote}
Before the new one: say the Mandarin for to be windy, then the Mandarin for to thunder.
\end{quote}

\subsection*{You'll want to know: liángkuai}
\emph{liángkuai} — \textquotedblleft{}pleasantly cool\textquotedblright{}.

\begin{quote}
\emph{qiūtiān hěn liángkuai} — Autumn is nice and cool.
\end{quote}

Its characters come later; for now it is a word for the ear.

\begin{cousinweb}
\emph{liángkuai} is built from \emph{liáng} 'cool' and \zh{快}, glossed 'to one's liking', both inherited from Middle Chinese, \emph{ljang} and \emph{khwaejH} (Guangyun, in Baxter's transcription): cool and pleasant.
\end{cousinweb}

\begin{grammarlens}[title={What you've built}]
One more word for giving reasons.
\end{grammarlens}

\subsection*{Your turn}
\begin{itemize}
\raggedright
  \item \emph{Say it:} \emph{liángkuai}
  \item \emph{Say it:} \emph{liángkuai}, once more
  \item \emph{Say it:} \emph{qiūtiān hěn liángkuai} — Autumn is nice and cool.
  \item \emph{Recall:} say the Mandarin for to be windy, then the Mandarin for to thunder, then say \emph{liángkuai} again
\end{itemize}

\begin{tcolorbox}[breakable,colback=teal!4,colframe=teal!35!black,title={Before you move on}]
What does liángkuai mean? (\textquotedblleft{}Pleasantly cool\textquotedblright{}.) Say it once more.
\end{tcolorbox}
